\documentclass[../../../include/open-logic-section]{subfiles}

\begin{document}

\olfileid[id]{sth}{spine}{Valphabasic}
\olsection{Sifat-Sifat Dasar Tahap}

Untuk menonjolkan arti penting definisi $V_\alpha$ bagi landasan teori, kita
akan menyajikan beberapa hasil dasar mengenainya. Kita mulai dengan sebuah
definisi:\footnote{Tidak ada istilah baku untuk ``poten''; nama inilah yang
digunakan oleh \citet{ButtonLT1}.}
\begin{defn}
	Himpunan $A$ disebut \emph{poten} jika dan hanya jika
	$\forall x((\exists y \in A)x \subseteq y \lif x \in A)$.
\end{defn}
\begin{lem}\ollabel{Valphabasicprops}
Untuk setiap ordinal $\alpha$:
\begin{enumerate}
	\item\ollabel{Valphatrans} Setiap $V_\alpha$ bersifat transitif.
	\item\ollabel{Valphapotent} Setiap $V_\alpha$ bersifat poten.
	\item\ollabel{Valphacum} Jika $\gamma \in \alpha$, maka $V_\gamma
	\in V_\alpha$ (dan karena itu juga $V_\gamma \subseteq V_\alpha$
	berdasarkan \olref{Valphatrans}).
\end{enumerate}
\end{lem}

\begin{proof}
Kita membuktikannya dengan induksi transfinit (serentak). Untuk langkah
induksi, andaikan \olref{Valphatrans}--\olref{Valphacum} berlaku bagi setiap
ordinal $\beta < \alpha$.

Kasus $\alpha = \emptyset$ bersifat langsung.

Andaikan $\alpha = \ordsucc{\beta}$. Untuk menunjukkan
\olref{Valphacum}, jika $\gamma \in \alpha$, maka
$V_\gamma \subseteq V_\beta$ menurut hipotesis, sehingga
$V_\gamma \in \Pow{V_\beta} = V_\alpha$. Untuk menunjukkan
\olref{Valphapotent}, andaikan $A \subseteq B \in V_\alpha$, yakni
$A \subseteq B \subseteq V_\beta$; maka $A \subseteq V_\beta$, sehingga
$A \in V_\alpha$. Untuk menunjukkan \olref{Valphatrans}, perhatikan bahwa
jika $x \in A \in V_\alpha$, maka $A \subseteq V_\beta$, sehingga
$x \in V_\beta$, lalu $x \subseteq V_\beta$ karena $V_\beta$ transitif
menurut hipotesis, dan dengan demikian $x \in V_\alpha$.

Andaikan $\alpha$ merupakan ordinal limit. Untuk menunjukkan
\olref{Valphacum}, jika $\gamma \in \alpha$, maka
$\gamma \in \ordsucc{\gamma} \in \alpha$, sehingga
$V_\gamma \in V_{\ordsucc{\gamma}}$ menurut asumsi, dan karenanya
$V_\gamma \in \bigcup_{\beta \in \alpha} V_\beta = V_\alpha$. Untuk
menunjukkan \olref{Valphatrans} dan \olref{Valphapotent}, cukup perhatikan
bahwa gabungan himpunan-himpunan transitif (berturut-turut, poten) bersifat
transitif (berturut-turut, poten).
\end{proof}

\begin{lem}\ollabel{Valphanotref}
Untuk setiap ordinal $\alpha$, $V_\alpha \notin V_\alpha$.
\end{lem}

\begin{proof}
Dengan induksi transfinit. Jelas bahwa
$V_\emptyset \notin V_\emptyset$.

Jika $V_{\ordsucc{\alpha}} \in V_{\ordsucc{\alpha}} = \Pow{V_\alpha}$,
maka $V_{\ordsucc{\alpha}} \subseteq V_\alpha$; dan karena
$V_\alpha \in V_{\ordsucc{\alpha}}$ berdasarkan
\olref{Valphabasicprops}, kita memperoleh $V_\alpha \in V_\alpha$.
Dengan mengambil kontraposisinya: jika $V_\alpha \notin V_\alpha$, maka
$V_{\ordsucc{\alpha}} \notin V_{\ordsucc{\alpha}}$.

Jika $\alpha$ merupakan limit dan
$V_\alpha \in V_\alpha = \bigcup_{\beta \in \alpha}V_\beta$, maka
$V_\alpha \in V_\beta$ untuk suatu $\beta \in \alpha$; tetapi kemudian
$V_\beta \in V_\alpha$ juga, sehingga $V_\beta \in V_\beta$ berdasarkan
\olref{Valphabasicprops} (dua kali). Dengan mengambil kontraposisinya, jika
$V_\beta \notin V_\beta$ untuk semua $\beta \in \alpha$, maka
$V_\alpha \notin V_\alpha$.
\end{proof}

\begin{cor}
Untuk sembarang ordinal $\alpha, \beta$: $\alpha \in \beta$ jika dan hanya
jika $V_\alpha \in V_\beta$.
\end{cor}

\begin{proof}
\Olref{Valphabasicprops} memberikan satu arah. Sebaliknya, andaikan
$V_\alpha \in V_\beta$. Maka $\alpha \neq \beta$ berdasarkan
\olref{Valphanotref}; dan $\beta \notin \alpha$, sebab jika tidak, kita akan
memperoleh $V_\beta \in V_\alpha$, dan kemudian
$V_\beta \in V_\beta$ berdasarkan \olref{Valphabasicprops} (dua kali),
bertentangan dengan \olref{Valphanotref}. Jadi, $\alpha \in \beta$
berdasarkan Trikotomi.
\end{proof}
\noindent
Semua ini memungkinkan kita memandang setiap $V_\alpha$ sebagai tahap
ke-$\alpha$ dalam hierarki. Berikut alasannya.

$V_\alpha$ kita tentu dapat dipandang terbentuk melalui suatu proses
\emph{iteratif}, sebab penggunaan ordinal oleh kita melacak gagasan iterasi.
Selain itu, jika suatu tahap terbentuk sebelum tahap lain, yakni
$V_\beta \in V_\alpha$, yakni $\beta \in \alpha$, maka proses pembentukan
kita bersifat \emph{kumulatif}, karena $V_\beta \subseteq V_\alpha$.
Akhirnya, kita memang membentuk \emph{semua} koleksi himpunan yang mungkin
yang tersedia pada suatu tahap sebelumnya, sebab setiap tahap penerus
$V_{\ordsucc{\alpha}}$ merupakan himpunan kuasa pendahulunya $V_\alpha$.

Singkatnya: dengan $\ZFminus$, kita \emph{hampir} selesai merumuskan visi
kita tentang hierarki kumulatif-iteratif himpunan. (Meskipun, tentu saja, kita
masih perlu membenarkan Penggantian.)

\end{document}
